\chapter{Testing}
\label{ch:testing}

\section{Testing Strategy}
The system was validated using a layered testing strategy that
combined static analysis, unit tests on the Python backend, end-to-end
manual UI walkthroughs, and a fifteen-participant user study. The
strategy was designed to catch defects at the lowest cost: failures
of the prompt schema were caught by Pydantic at the request boundary;
LLM-side regressions were caught by golden-output comparison; and
multimodal pipeline failures were detected through deliberate
fault-injection (e.g.\ rate-limited image responses).

\section{Unit Testing}
Unit tests cover the deterministic portions of the backend:
context-extraction projections, the five-phase finite-state machine,
the prompt-builder template engine and the SQLite persistence layer.
Each test is parameterised across all four narrative phases. The
test suite is executed with \texttt{pytest} and runs in under three
seconds on the development machine.

\begin{table}[!t]
\caption{Unit-test coverage by module.}
\label{tab:unittests}
\centering
\renewcommand{\arraystretch}{1.15}
\begin{tabular}{@{}lcc@{}}
\toprule
\textbf{Module}                & \textbf{Tests} & \textbf{Line cov.} \\
\midrule
\texttt{user\_input}            & 14 & 92\% \\
\texttt{context\_extraction}    & 22 & 89\% \\
\texttt{knowledge\_memory}      & 11 & 84\% \\
\texttt{prompt\_engineering}    & 18 & 91\% \\
\texttt{story\_flow\_manager}   & 16 & 88\% \\
\texttt{output\_formatter}      &  9 & 95\% \\
\midrule
\textbf{Overall}                & \textbf{90} & \textbf{89.8\%} \\
\bottomrule
\end{tabular}
\end{table}

\section{Integration Testing}
End-to-end integration tests exercise the full session lifecycle:
\begin{enumerate}
    \item \texttt{POST /api/session/new} -- create session.
    \item \texttt{GET  /api/questions/all} -- fetch all questions.
    \item \texttt{POST /api/response} (loop) -- submit answers.
    \item \texttt{POST /api/story/generate} -- generate first segment.
    \item \texttt{POST /api/story/continue} -- continue from choice.
    \item \texttt{GET  /api/story/full/\{id\}} -- retrieve full story.
    \item \texttt{DELETE /api/session/\{id\}} -- end session.
\end{enumerate}
All seven steps are exercised once per provider (Ollama, Groq,
Hugging Face, OpenAI). Integration tests use \texttt{httpx.AsyncClient}
hitting an in-process FastAPI instance.

\section{API Testing}
Each REST endpoint was validated for: (i) correct status code on
success and failure, (ii) Pydantic-schema rejection of malformed
payloads, (iii) idempotency of session reads, (iv) graceful
handling of unknown session IDs (404 rather than 500). A summary
of API test outcomes appears in Table~\ref{tab:apitests}.

\begin{table}[!t]
\caption{API testing outcomes.}
\label{tab:apitests}
\centering
\renewcommand{\arraystretch}{1.15}
\small
\begin{tabular}{@{}llcc@{}}
\toprule
\textbf{Endpoint} & \textbf{Tests} & \textbf{Pass} & \textbf{Notes} \\
\midrule
\texttt{/api/status}                  & 4  & 4  & Provider availability \\
\texttt{/api/session/new}             & 3  & 3  & UUID format \\
\texttt{/api/questions/all}           & 2  & 2  & Phase order \\
\texttt{/api/response}                & 8  & 8  & Including custom\_q\_N \\
\texttt{/api/story/generate}          & 6  & 6  & All 4 providers \\
\texttt{/api/story/continue}          & 6  & 6  & With/without choice \\
\texttt{/api/settings}                & 5  & 5  & Hot-swap validated \\
\texttt{/api/session/\{id\}/summary}  & 3  & 3  & Word/segment counts \\
\bottomrule
\end{tabular}
\end{table}

\section{UI / End-to-End Testing}
Manual end-to-end walkthroughs were performed against the four
browsers used during development (Chrome 124, Safari 17, Firefox 124,
Edge 124). The walkthrough covers: question budget selection,
question form filling with custom answers, custom-question authoring,
story generation, image rendering, narration playback, ambient music
toggle, Story Map modal, Character Graph modal and the admin panel.

\section{Multimodal Robustness Testing}

\subsection{Image-API Fault Injection}
To validate the throttle queue's robustness, we deliberately fired
five concurrent image-generation requests across the prompt-length
range 80--380 characters. Without the queue, the failure probability
rose from $\approx$40\% at 80 chars to $\approx$98\% at 380 chars; with
the serial queue and 1.2\,s spacing, failures dropped to below 2\%
across the same range. These results are quantified in
Chapter~\ref{ch:results}.

\subsection{Audio Autoplay Policy}
The browser autoplay policy blocks programmatic
\texttt{Audio.play()} unless a recent user gesture has occurred.
We verified that the ambient-music toggle satisfies this requirement by
delaying the first \texttt{play()} call until the user clicks the
toggle on the actual story view; once unblocked, subsequent plays
work for the session lifetime.

\subsection{Cross-Origin Audio}
A defect in an early implementation set
\texttt{audio.crossOrigin = 'anonymous'} which caused the audio
load to fail because SoundHelix does not return
\texttt{Access-Control-Allow-Origin} headers. The fix --- removing
the cross-origin attribute --- was verified across all four
browsers.

\section{Consistency-Checker Evaluation}
The active consistency checker was tested against 50 ground-truth
annotated story segments. Each segment was checked for five issue
classes (character contradictions, plot-thread abandonment, pacing
issues, setting inconsistencies, tone shifts). Detection rates
ranged from 68\% (tone shifts) to 87\% (character contradictions);
detailed numbers appear in Chapter~\ref{ch:results}.

\section{Performance Testing}
End-to-end latency was benchmarked across all four providers. Groq
LLaMA-3.1-8B is the fastest at 1.8\,s mean initial segment, while
local Ollama runs in 12.4\,s on the development MacBook. The
detailed breakdown appears in Figure~\ref{fig:latency-results}
of Chapter~\ref{ch:results}.

\section{Acceptance Criteria}
The system was deemed to meet acceptance when:
\begin{itemize}
    \item All seven integration steps complete on at least three of
    four providers without manual intervention.
    \item Mean Likert score across the four narrative dimensions
    exceeds 3.8 / 5.
    \item Image-generation success rate exceeds 95\% after retries.
    \item Story-Map and Character-Graph modals render in under
    500\,ms for stories of up to ten segments.
    \item No critical bug remains open after the final user-study
    debrief.
\end{itemize}
All acceptance criteria were met before submission.
